% Generated by scripts/materialize_unified_paper.py; do not edit.
\section{Reproducibility and artifact publication}\label{fin:sec:repro}

\subsection{Two verification tiers}

The repository package~\cite{Repository2026} separates lightweight tracked
checks from large proof replay.  At the date of this draft, public GitHub and
archival release publication remain pending; the following is the release
contract, not a claim that an anonymous download already exists.

\paragraph{Tier 1: source and semantics.}
Python \(3.11.15\) and the locked environment reproduce graph verification,
unit tests, and exact-seven frozen-state checks.  From the repository root,
the lightweight entry point is
\begingroup\scriptsize
\begin{verbatim}
bash scripts/bootstrap.sh
.venv/bin/python reproduce.py quick
\end{verbatim}
\endgroup
The proof-critical Python packages are pinned to
\texttt{python-flint==0.9.0}, \texttt{mpmath==1.4.1},
\texttt{python-sat==1.9.dev13}, and \texttt{six==1.17.0}.  The finite route
uses PySAT; the other entries are repository-wide proof dependencies.

\paragraph{Tier 2: large proof replay.}
After downloading the release assets into one directory, the content manifest
is checked before decompression:
\begingroup\small
\begin{verbatim}
.venv/bin/python scripts/verify_artifacts.py \
  --directory /path/to/release-assets --exact
\end{verbatim}
\endgroup
The exact mode also rejects unmanifested files, so a stale or mixed download
directory cannot masquerade as the immutable release payload.
Build \texttt{drat-trim} from the pinned source commit with
\path{scripts/build_drat_trim.sh}; a platform-specific executable hash is
recorded but not used as a cross-platform acceptance condition.  The one
integrated command
\begin{verbatim}
.venv/bin/python reproduce.py finite-heavy \
  --artifact-dir /path/to/release-assets \
  --drat-trim /absolute/path/to/drat-trim
\end{verbatim}
materializes a temporary overlay containing tracked semantic metadata and the
twelve release payloads.  It reconstructs all six formulas, replays all six
proofs, and audits the exact-seven state.  A proof is accepted only when the
checker exits zero and prints \texttt{s VERIFIED}.

\subsection{Why large proofs are release assets}

The six exact-six compressed files occupy about \(1.26\) GB, while the three
raw DRAT traces occupy about \(6.39\) GB.  Adding the smaller budget-five
proofs to ordinary Git history would still make cloning needlessly expensive.
The repository therefore tracks code, matrices, cut banks, semantic records,
tests, the paper, and a content-addressed manifest.  CNF and DRAT payloads are
to be published as immutable release or archival assets (for example, a
GitHub Release mirrored to Zenodo), keyed by SHA-256 and byte count.  This design
keeps ordinary clones reviewable while preserving bit-for-bit replay.

The release is complete only if it includes both the three budget-five proof
pairs used by \cref{fin:prop:budget5} and the three exact-six proof pairs used by
\cref{fin:prop:unsat}.  Their identities are listed in
\cref{fin:tab:budget5-artifacts,fin:tab:budget6-artifacts}.  Incomplete, corrupt, and
exploratory traces are explicitly excluded.

\subsection{Reproducibility contract}

A reproduction should report four logically distinct outcomes:
\begin{enumerate}[label=(\arabic*)]
\item the seed and artifact hashes match;
\item semantic reconstruction passes;
\item each required DRAT replay is verified; and
\item the exact-seven endpoint remains \textsc{unknown}, unless a new,
  separately reviewed certificate supersedes it.
\end{enumerate}
Passing only a solver run, only a hash check, or only unit tests does not by
itself reproduce the full theorem.
